%======================================================================
% Preface: how to use these notes
%======================================================================
\chapter*{How to Use These Notes}\label{ch:preface}
\addcontentsline{toc}{chapter}{How to Use These Notes}
\markboth{How to Use These Notes}{How to Use These Notes}

\section*{Who these notes are for}

These notes are written for someone who wants to understand machine learning
(ML) and deep learning (DL) properly, but who currently knows little or no
mathematics beyond everyday arithmetic: adding, subtracting, multiplying and
dividing numbers. If you have ever opened a machine learning paper, a textbook
or even a tutorial and been stopped by a wall of symbols such as
$\nabla_{\vtheta} L$, $\mW^\top\vx$, $\E[X]$ or $\sum_{i=1}^n$, these notes are
meant to take you from that point to the point where those symbols feel like
old friends.

Nothing is assumed. Every concept is introduced from first principles, with a
reason for why it exists, a picture of what it means, concrete numerical
examples, and only then the general rule. When we say ``from absolute
beginner'', we mean it: Chapter~\ref{ch:numbers} starts with what a number
is.

At the same time, the destination is ambitious. By the last chapter you will
have met every piece of mathematics that a standard first course in machine
learning or deep learning takes for granted: linear algebra (vectors, matrices,
eigenvalues, the singular value decomposition), calculus (derivatives,
gradients, the chain rule, Taylor approximations), optimization (gradient
descent and its relatives), probability and statistics (distributions,
expectation, Bayes' theorem, maximum likelihood, the bias--variance tradeoff),
and the specific tools of deep learning (softmax, cross-entropy, matrix
calculus and backpropagation).

\section*{What these notes are not}

They are not a course on programming, and they are not a complete treatment of
any branch of mathematics. Where a topic is deep (for example, the theory of
limits or of measure-theoretic probability), we give the intuition that a
practitioner needs and point you to where the rigorous story is told. The
final chapter, \emph{Where to Go Next}, suggests what to read after these
notes.

\section*{How the notes are organised}

The notes are divided into twelve chapters that build on one another. The
diagram below shows which chapters depend on which. The three main
``strands'' are linear algebra (left), calculus and optimization (middle) and
probability and statistics (right); they come together in
Chapter~\ref{ch:dl}, where we assemble the mathematics of a neural network.

\begin{center}
\begin{tikzpicture}[
  node distance=7mm and 6mm,
  chap/.style={draw,rounded corners=3pt,thick,align=center,
    minimum width=33mm,minimum height=10mm,font=\small,fill=white},
  found/.style={chap,draw=MLremark,fill=MLremark!8},
  linalg/.style={chap,draw=MLdef,fill=MLdef!8},
  calc/.style={chap,draw=MLexample,fill=MLexample!8},
  prob/.style={chap,draw=MLthm,fill=MLthm!8},
  dl/.style={chap,draw=MLml,fill=MLml!10},
  dep/.style={-{Stealth[length=2mm]},thick,gray!70},
]
  \node[found] (c1) {\textbf{1} Numbers \&\\ algebra};
  \node[found,below=of c1] (c2) {\textbf{2} Functions \&\\ graphs};
  \node[found,left=of c2] (c3) {\textbf{3} Geometry \&\\ trigonometry};
  \node[linalg,below=of c3] (c4) {\textbf{4} Linear algebra I:\\ vectors \& matrices};
  \node[calc,below=of c2] (c6) {\textbf{6} Calculus I:\\ one variable};
  \node[prob,right=of c6] (c9) {\textbf{9} Probability};
  \node[linalg,below=of c4] (c5) {\textbf{5} Linear algebra II:\\ core ML topics};
  \node[calc,below=of c6] (c7) {\textbf{7} Multivariable\\ calculus};
  \node[prob,below=of c9] (c10) {\textbf{10} Statistics\\ for ML};
  \node[calc,below=of c7] (c8) {\textbf{8} Optimization};
  \node[dl,below=of c8] (c11) {\textbf{11} Topics for\\ deep learning};
  \node[found,below=of c11] (c12) {\textbf{12} Where to go next};

  \draw[dep] (c1) -- (c2);
  \draw[dep] (c2) -- (c3);
  \draw[dep] (c3) -- (c4);
  \draw[dep] (c4) -- (c5);
  \draw[dep] (c2) -- (c6);
  \draw[dep] (c6) -- (c7);
  \draw[dep] (c5) -- (c7);
  \draw[dep] (c6) -- (c9);
  \draw[dep] (c9) -- (c10);
  \draw[dep] (c7) -- (c8);
  \draw[dep] (c5.south) |- (c11.west);
  \draw[dep] (c8) -- (c11);
  \draw[dep] (c10.south) |- (c11.east);
  \draw[dep] (c11) -- (c12);
\end{tikzpicture}
\end{center}

If you are a complete beginner, read the chapters in order. If you already
feel comfortable with some material (say, you remember high-school algebra),
skim the corresponding chapter, try its exercises, and move on as soon as they
feel easy. Each chapter opens with a list of what you will learn, which you can
use as a checklist.

\section*{The boxes}

The text is interspersed with coloured boxes. Each colour always means the
same thing, so that you can find your way around quickly.

\begin{tcolorbox}[mlbox=MLmotivation]
\textbf{\color{MLmotivation}Motivation.} Opens every chapter and many sections:
\emph{why} do we need this idea, and where does it show up in machine learning?
\end{tcolorbox}
\begin{tcolorbox}[mlbox=MLdef]
\textbf{\color{MLdef}Definition.} The precise meaning of a new word or symbol.
Definitions are the vocabulary of mathematics; it pays to read them twice.
\end{tcolorbox}
\begin{tcolorbox}[mlbox=MLthm]
\textbf{\color{MLthm}Theorem / Proposition / Corollary.} A fact that follows from
the definitions. Often we explain why it is true, sometimes with a short proof.
\end{tcolorbox}
\begin{tcolorbox}[mlbox=MLexample,colback=MLexample!3!white]
\textbf{\color{MLexample}Example.} A worked example with concrete numbers. Try to
solve it yourself before reading the \emph{Solution}.
\end{tcolorbox}
\begin{tcolorbox}[mlbox=MLintuition,colback=MLintuition!7!white]
\textbf{\color{MLintuition}Intuition / Picture this.} An informal explanation, an
analogy, or a mental picture. These are often the most important boxes of all.
\end{tcolorbox}
\begin{tcolorbox}[mlbox=MLml]
\textbf{\color{MLml}Why this matters for ML.} How the idea is used in machine
learning or deep learning.
\end{tcolorbox}
\begin{tcolorbox}[mlbox=MLpitfall]
\textbf{\color{MLpitfall}Common pitfall.} A mistake that many people make, and how
to avoid it.
\end{tcolorbox}
\begin{tcolorbox}[mlbox=MLnotation]
\textbf{\color{MLnotation}Notation.} How a new symbol is written and read aloud.
\end{tcolorbox}
\begin{tcolorbox}[mlbox=MLremark,colback=MLremark!4!white]
\textbf{\color{MLremark}Remark.} A side comment: a subtlety, a generalisation, or a
historical note. Safe to skip on a first reading.
\end{tcolorbox}
\begin{tcolorbox}[mlbox=MLalgorithm]
\textbf{\color{MLalgorithm}Algorithm.} A recipe, written as a list of steps, that a
computer (or you, with a pencil) can follow.
\end{tcolorbox}
\begin{tcolorbox}[mlbox=MLexercise,colback=MLexercise!3!white]
\textbf{\color{MLexercise}Exercise.} A problem for you to solve. Exercises appear at
the end of most sections. Every exercise has a solution in
Appendix~\ref{app:solutions}; the small link in the bottom-right corner of
each exercise box takes you straight to it, and the exercise number in the
appendix takes you back.
\end{tcolorbox}

New technical terms are printed in \textbf{\color{MLterm}bold colour} where
they are first explained, and every such term is listed in the index at the
end of the notes.

\section*{How to study mathematics}

Mathematics is not a spectator sport. The following habits make an enormous
difference.

\begin{itemize}
  \item \textbf{Read slowly, with pencil and paper.} A page of mathematics can take
    much longer to read than a page of a novel. That is normal. When the text
    says ``it follows that'', check that it does.
  \item \textbf{Work through every example yourself.} Cover the solution, try the
    problem, then compare.
  \item \textbf{Do the exercises.} Understanding comes from doing. Start with the
    first few exercises of each section; they are the easiest.
  \item \textbf{Say the symbols out loud.} Every formula is a sentence. The
    \emph{Notation} boxes tell you how to read new symbols, and
    Appendix~\ref{app:notation} collects them all.
  \item \textbf{Expect confusion, and push through it.} Being confused is not a
    sign that you are bad at mathematics; it is what learning mathematics feels
    like. Come back to a difficult section after a day or two and it will often
    have become clear.
  \item \textbf{Connect to code if you can.} If you program, implement small
    examples (for instance in Python with NumPy). Seeing a formula produce the
    same number as your hand calculation is a powerful check.
\end{itemize}

\section*{Reference material at the end}

Appendix~\ref{app:notation} lists all notation and the Greek alphabet,
Appendix~\ref{app:formulas} is a compact formula sheet for revision, and
Appendix~\ref{app:solutions} contains solutions to all exercises.
